\chapter{Prezentarea proiectului}

\section{Precizări asupra conținutului și a modului de organizare}

Împreună cu capitolul următor reprezintă cca. 70\% din lucrare. 

Titlul acestui capitol nu este unul impus și nici nu corespunde neapărat unui singur capitol. 
Titlul indică mai degrabă o parte (importantă și centrală, de altfel) a lucrării, în care se prezintă 
ceea ce s-a realizat efectiv: contribuțiile autorului. Organizarea acestei părți este dependentă și 
specifică fiecărei lucrări în parte și este stabilită de către fiecare autor după cum i se pare mai 
potrivit pentru tema lui. Ea poate cuprinde prezentarea unor concepte teoretice (unelte sau tehnici
matematice folosite în lucrare, prezentarea sau introducerea unor concepte teoretice etc.),  o 
analiză a diferitelor metode/algoritmi/tehnologii etc. luate în considerare sau dezvoltate de către 
autor, o prezentare a unui design (mai mult sau mai puțin detaliat) sau chiar detalii a unei 
eventuale implementări/prototip, dacă e cazul.

Trebuie remarcat însă faptul că această parte reprezintă contribuția personală a autorului, și în nici 
un caz ea nu poate fi sinteza unor texte preluate din alte surse. Prin urmare, orice informații sunt 
prezentate aici, ele trebuie să corespundă cel puțin unei interpretări/analize critice personale a 
autorului, dacă nu chiar unor idei originale ale acestuia. 


\subsection{Subsecțiune}

